\scanpage{82}
% archive foliation: 81
% lines 1--3
des $\Phi_i\in\operatorname{Homext}(T,\pi)$ canoniques,
$i\in I$. Le système
\deleted{$\unsure{\Sigma(X/I/k)|\Sigma(X/I/x/k)}$}\inserted{$\Sigma(X/I/x/k)$}
$=(T,I,\pi,(\Phi_i))$
satisfait aux conditions envisagées. \deleted{\unsure{Pour}}
% lines 4--7
\deleted{\unsure{Changeons}} \unsure{On}
\inserted{Au lieu de} prendre $x$, il suffit
de prendre un $\unsure{\xi|\zeta}$, foncteur fibre sur
la catégorie des rev. étales de $X\setminus I$
($\pi$ étant le groupe des autom. de $\unsure{\xi|\zeta}$\ldots)
% lines 8--11
Alors $\Sigma$ dépend fonctoriellement de
$(X/I/\unsure{\xi|\zeta}/k)$. Prenons $X/I$, $k$ fixes,
$\unsure{\xi|\zeta}$ variable, on voit que $\Sigma$ est déterminé
``mod autom. intérieurs'' i.e. \deleted{\unsure{on a}}
% lines 12--17
définit un torseur \deleted{sous}\inserted{à droite} (non pas
sous $\widetilde A_\nu$ mais sous son quotient
$\widetilde A_\nu/\pi_\nu=)\ \widetilde\Gamma_\nu$,
\deleted{le dernier torseur}\inserted{soit $P(X/I/k)$},
dont le commutant est le groupe
des ``autom. ext.'' de \deleted{$\unsure{(T,I,\pi)}$}
$\Sigma(X/I/\unsure{x|\xi}/k)\ldots$
% lines 18--20
Supposons que $X/I/\unsure{k}$ provienne
par extension de la base
\deleted{$\unsure{\operatorname{Spec}k\longrightarrow S_0}$}\inserted{$k_0\subset k$}
\deleted{d'une situation}
